\documentclass[aspectratio=169]{beamer}

% --------------------------------------------------
% Theme and appearance
% --------------------------------------------------

\usetheme{Madrid}
\usecolortheme{default}

\setbeamertemplate{navigation symbols}{}
\setbeamertemplate{caption}[numbered]

% --------------------------------------------------
% Packages
% --------------------------------------------------

\usepackage[T1]{fontenc}
\usepackage{lmodern}
\usepackage{amsmath}
\usepackage{amssymb}
\usepackage{booktabs}
\usepackage{array}
\usepackage{graphicx}
\usepackage{multicol}
\usepackage{xcolor}

% --------------------------------------------------
% Title information
% --------------------------------------------------

\title{Barcoding Score Test}
\subtitle{An Interpretable Sliding-Window Detection Concept}
\author{Jonathan Ma}
\institute{Johns Hopkins University}
\date{\today}

% --------------------------------------------------
% Helpful commands
% --------------------------------------------------

\newcommand{\feature}[1]{\textbf{#1}}
\newcommand{\code}[1]{\texttt{#1}}

\begin{document}

% ==================================================
% Title
% ==================================================

\begin{frame}
    \titlepage
\end{frame}

% ==================================================
% Overview
% ==================================================

\begin{frame}{Barcoding Score Test}

\textbf{Objective}
Develop an interpretable, classical image-analysis detector that assigns a
continuous barcoding score to each horizontal position of a processed OCT
B-scan.
\vspace{0.4cm}

\textbf{General workflow}
\[
\text{E2E volume}
\rightarrow
\text{Preprocess \& Denoise}
\rightarrow
\text{Feature Engineering}
\rightarrow
\text{Barcoding Score}
\]

\vspace{0.4cm}

The initial detector evaluates two forms of evidence:

\begin{itemize}
    \item \textbf{Structural evidence}: vertical organization and repeated
    spatial texture.
    \item \textbf{Hypertransmission evidence}: brightness and persistence
    through depth.
\end{itemize}

\vspace{0.2cm}

This first run is intended as a proof of concept. No final clinical threshold,
feature weights, or parameter values have yet been tuned.

\end{frame}

% ==================================================
% Preprocessing
% ==================================================

\section{Preprocessing}

\begin{frame}{Preprocessing Pipeline}

For each selected E2E volume:

\begin{enumerate}
    \item Load, then Select the center B-scan using
    \[
    i_{\mathrm{center}}
    =
    \left\lfloor
    \frac{N_{\mathrm{B\text{-}scans}}}{2}
    \right\rfloor.
    \]
    \item Retrieve the BM boundary and interpolate missing boundary coordinates.
    \item Flatten the image so BM lies on one horizontal reference row.
    \item Crop a fixed region extending 150 pixels below BM.
    \item Apply whole-ROI z-score normalization.
    \item Apply mild anisotropic Gaussian denoising.
\end{enumerate}

\vspace{0.2cm}

The resulting detector input has approximate dimensions $150(h) \times 512(w)$

\end{frame}

\begin{frame}{Whole-ROI Z-Score Normalization}

After flattening, the entire  sub-BM region is normalized using
\[
Z(z,x)
=
\frac{
I_{\mathrm{ROI}}(z,x)-\mu_{\mathrm{ROI}}
}{
\sigma_{\mathrm{ROI}}+\varepsilon
},
\]
where: \(I_{\mathrm{ROI}}(z,x)\) is the intensity and \(\varepsilon=10^{-8}\) prevents division by zero.

\vspace{0.3cm}

For the first test scan:

\[
\mu_{\mathrm{ROI}} \approx 60.98,
\qquad
\sigma_{\mathrm{ROI}} \approx 44.73.
\]

\vspace{0.3cm}

Normalization is performed once over the whole ROI; the sliding-window features
do not renormalize each window.

\end{frame}

\begin{frame}{Gaussian (Pseeudo-ImageJ) Denoising }

A two-dimensional anisotropic Gaussian filter is applied after normalization:

\[
I_{\sigma_z,\sigma_x}
=
G_{\sigma_z,\sigma_x}
*
Z,
\]

where:

\begin{itemize}
    \item \(G_{\sigma_z,\sigma_x}\) is a Gaussian smoothing kernel,
    \item \(*\) denotes convolution,
    \item \(\sigma_z\) controls smoothing along depth, 1.0 first used
    \item \(\sigma_x\) controls smoothing horizontally, 0.5 first used
\end{itemize}

\vspace{0.3cm}

\textbf{Reason for asymmetric smoothing}

\begin{itemize}
    \item More smoothing is applied through depth to reduce OCT speckle.
    \item Less horizontal smoothing is applied to avoid erasing narrow
    vertical stripe boundaries.
    \item The observed denoising effect was intentionally subtle rather than
    strongly blurring the scan.
\end{itemize}

\end{frame}

\begin{frame}{Preprocessing Visualization}

\begin{center}
\textbf{preprocess\_bscan(volume, index, config)}
\end{center}

\end{frame}

\begin{frame}{Sliding-Window Configuration}

The detector uses overlapping vertical windows that span the complete
150-pixel depth.

For a window centered at horizontal position \(x\),
\[
W_x
=
I
\left[
0:149,\,
x-r:x+r
\right],
\]
where the initial window radius is r=10. Therefore,

\[
\text{window width}=2r+1=21 \text{ pixels}.
\]

\vspace{0.3cm}

The initial stride is a single pixel, so the windows are:
\[
(0,\ldots,20),
\quad
(1,\ldots,21),
\quad
(2,\ldots,22),
\quad \ldots
\]

Each window produces one feature vector assigned to its center position.
Reflection padding is used at the horizontal image boundaries. The first and
last 10 score positions are later excluded from detection to reduce edge
artifacts.

\end{frame}

% ==================================================
% Structural features
% ==================================================

\section{Structural Feature Building}

\begin{frame}{Structural Feature 1: Verticality}

Verticality is calculated using a \textbf{structure tensor}, which summarizes
the direction of local image-intensity changes.

Let
\[
I_x=\frac{\partial I}{\partial x},
\qquad
I_z=\frac{\partial I}{\partial z},
\]
where:
\begin{itemize}
    \item \(I_x\) measures horizontal brightness change,
    \item \(I_z\) measures brightness change through depth.
\end{itemize}
The local structure tensor is
\[
J
=
\begin{bmatrix}
\langle I_x^2\rangle & \langle I_xI_z\rangle \\
\langle I_xI_z\rangle & \langle I_z^2\rangle
\end{bmatrix}.
\]

Its eigenvalues satisfy $\lambda_1\geq \lambda_2\geq 0$. We calculate an anisotropy score as
\[
A
=
\frac{
\lambda_1-\lambda_2
}{
\lambda_1+\lambda_2+\varepsilon
}.
\]
\end{frame}

\begin{frame}{Verticality: Interpretation}

The anisotropy score describes how strongly the local texture prefers one
direction, with greater values indicating strong orientation
A vertical line usually produces strong horizontal intensity changes.
Therefore, the verticality map combines:

\begin{itemize}
    \item the strength of directional organization;
    \item agreement with the orientation expected from a vertical structure.
\end{itemize}

The window-level verticality score is the mean verticality across all pixels in
the local 150-by-21 window.

\vspace{0.3cm}

\textbf{Observed first-run behavior}

\begin{itemize}
    \item Verticality rose around visibly prominent vertical transmission
    columns.
    \item It was one of the strongest and most spatially interpretable
    features.
    \item Strong edge responses were also observed, motivating edge exclusion.
\end{itemize}

\end{frame}

\begin{frame}{Structural Feature 2: Depth Continuity}

Depth continuity measures whether the horizontal bright--dark pattern remains
aligned as the image is traversed downward.

For each pair of neighboring depth rows, the detector calculates their
horizontal correlation:
\[
C_z
=
\operatorname{corr}
\left(
I(z,:),
I(z+1,:)
\right).
\]
The window continuity score is the mean of the valid row-pair correlations:
\[
C(x)
=
\frac{1}{N_{\mathrm{valid}}}
\sum_z C_z.
\]
\vspace{0.3cm}

Here:
\begin{itemize}
    \item \(I(z,:)\) is the horizontal intensity pattern at depth row \(z\),
    \item \(I(z+1,:)\) is the pattern one row deeper,
\end{itemize}
The initial depth lag is a single row.
\end{frame}

\begin{frame}{Depth Continuity: Interpretation}

A barcode-like region contains vertical structures, so bright and dark
horizontal positions should remain similar across successive depth rows.

\vspace{0.3cm}

\textbf{Interpretation}

\begin{itemize}
    \item High continuity: the same column pattern persists through depth.
    \item Low continuity: local brightness patterns change substantially
    between depth rows.
    \item Negative continuity: horizontal patterns tend to reverse.
\end{itemize}

\vspace{0.3cm}

\textbf{Observed first-run behavior}

\begin{itemize}
    \item Continuity aligned well with visually prominent vertical regions.
    \item It was substantially more stable than the initially tested
    periodicity feature.
    \item It was strongly correlated with verticality: (r=0.917)
    \item It was also strongly correlated with heterogeneity: (r=0.920)
\end{itemize}

These correlations motivated the grouped scoring formulation presented later.

\end{frame}

\begin{frame}{Structural Feature 3: Spatial Heterogeneity}

Spatial heterogeneity measures how much the local horizontal column intensity
varies inside a window.
First, each column is summarized across depth:
\[
c(u)
=
\frac{1}{D}
\sum_{z=0}^{D-1}
I(z,u),
\]
where:
\begin{itemize}
    \item \(u\) is a horizontal column inside the local window,
    \item \(D=150\) is the ROI depth,
    \item \(c(u)\) is the average intensity of that column.
\end{itemize}
The initial heterogeneity measure utilizes the standard deviation across the 21 column means.
\end{frame}

\begin{frame}{Spatial Heterogeneity: Interpretation}

\textbf{Interpretation}

\begin{itemize}
    \item Low heterogeneity: neighboring columns have similar average
    intensity.
    \item High heterogeneity: the window contains alternating brighter and
    darker columns or abrupt spatial changes.
\end{itemize}

This may help distinguish:

\begin{itemize}
    \item broad, relatively smooth hypertransmission;
    \item repeated or irregular vertical stripe structures.
\end{itemize}

\vspace{0.3cm}

\textbf{Observed first-run behavior}

Heterogeneity correlated highly with verticality (0.877) and continuity (0.920) so these three features are treated as related measurements of
\textbf{structural organization}, not as fully independent evidence.

\end{frame}

% ==================================================
% Hypertransmission features
% ==================================================

\section{Hypertransmission Feature Building}

\begin{frame}{Hypertransmission Feature 1: Depth Persistence}

Depth persistence measures how far candidate bright and vertically organized
signal extends through the 150-pixel ROI.

A candidate pixel is defined (loosely) using:

\begin{itemize}
    \item an intensity threshold at the 75th percentile of the processed scan;
    \item a minimum verticality requirement of 0.20.
\end{itemize}

For each column \(u\), define a binary candidate signal whose longest-run persistence is

\[
P(u)
=
\frac{
L_{\max}(u)
}{
D
},
\]

where:

\begin{itemize}
    \item \(L_{\max}(u)\) is the longest uninterrupted candidate run,
    \item \(D=150\) is the full ROI depth.
\end{itemize}

The sliding-window persistence score is the mean across the 21 columns.

\end{frame}

\begin{frame}{Depth Persistence: Interpretation}

\textbf{Interpretation}

\begin{itemize}
    \item \(P(u)\approx 0\): little bright, vertical signal persists through
    depth.
    \item Larger \(P(u)\): a longer continuous candidate signal is present.
    \item \(P(u)\approx 1\): candidate signal extends through nearly the entire
    ROI depth.
\end{itemize}

\vspace{0.3cm}

Persistence is intended to distinguish:

\begin{itemize}
    \item isolated bright pixels or shallow structures;
    \item sustained hypertransmission columns.
\end{itemize}

\vspace{0.3cm}

\textbf{Observed first-run behavior}

\begin{itemize}
    \item Persistence increased near the strongest visible transmission
    regions.
    \item It was strongly related to amplitude:
    \[
    r_{\mathrm{persistence,amplitude}}
    =
    0.823.
    \]
\end{itemize}

For this reason, persistence and amplitude form the
\textbf{hypertransmission feature group}.

\end{frame}

\begin{frame}{Hypertransmission Feature 2: Amplitude}

Amplitude measures the average normalized intensity within each 150-by-21
window:

\[
M(x)
=
\frac{1}{|W_x|}
\sum_{(z,u)\in W_x}
I(z,u).
\]

Here:

\begin{itemize}
    \item \(W_x\) is the window centered at horizontal position \(x\),
    \item \(|W_x|\) is the number of pixels in the window,
    \item \(I(z,u)\) is the normalized and denoised intensity.
\end{itemize}

\vspace{0.3cm}

\textbf{Interpretation}

\begin{itemize}
    \item Higher amplitude indicates a brighter local sub-BM region.
    \item This is consistent with stronger hypertransmission.
    \item Amplitude alone is not sufficient for barcoding because EA may also
    appear broadly bright.
\end{itemize}

Therefore, amplitude is combined with structural features rather than used as
a standalone classifier.

\end{frame}

% ==================================================
% Score construction
% ==================================================

\section{Barcoding Score Construction}

\begin{frame}{Robust Feature Standardization}

The raw feature signals have different numerical scales. Before weighting,
each feature is robustly standardized across the scan:
\[
\widetilde F_j(x)
=
\frac{
F_j(x)-\operatorname{median}(F_j)
}{
\operatorname{IQR}(F_j)+\varepsilon
}.
\]
Here:
\begin{itemize}
    \item \(F_j(x)\) is feature \(j\) at position \(x\),
    \item \(\operatorname{median}(F_j)\) is the median across the scan,
    \item \(\operatorname{IQR}(F_j)\) is the 75th percentile minus the 25th
    percentile,
    \item \(\varepsilon=10^{-8}\) prevents division by zero.
\end{itemize}
The standardized values are clipped to [-5,5] to limit the influence of extreme outliers.

\vspace{0.3cm}

This standardization makes the component magnitudes comparable before
constructing the weighted score.

\end{frame}

\begin{frame}{Individual Versus Grouped Scoring}

Two scoring modes were implemented.

\vspace{0.2cm}

\textbf{Individual scoring}

\[
S_{\mathrm{individual}}(x)
=
\sum_{j=1}^{5}
w_j\widetilde F_j(x).
\]

Each of the five features receives its own direct weight.

\vspace{0.3cm}

\textbf{Grouped scoring}

Highly correlated features are first combined into two interpretable groups:
\[
S_{\mathrm{structural}}(x)
=
w_V\widetilde V(x)
+
w_C\widetilde C(x)
+
w_H\widetilde H(x),
\]
\[
S_{\mathrm{hyperTD}}(x)
=
w_P\widetilde P(x)
+
w_A\widetilde A(x).
\]
The final score is
\[
S(x)
=
w_S S_{\mathrm{structural}}(x)
+
w_T S_{\mathrm{hyperTD}}(x).
\]
The grouped mode was selected as the initial default.

\end{frame}

\begin{frame}{Initial Individual Feature Weights}

For the direct five-feature score, the initial proof-of-concept weights are:

\begin{table}
\centering
\begin{tabular}{lc}
\toprule
Feature & Weight \\
\midrule
Verticality & 0.25 \\
Persistence & 0.25 \\
Continuity & 0.15 \\
Amplitude & 0.20 \\
Heterogeneity & 0.15 \\
\bottomrule
\end{tabular}
\end{table}

These weights sum to one.

\vspace{0.3cm}

They are not learned or clinically validated. They were chosen to:

\begin{itemize}
    \item give substantial weight to vertical geometry;
    \item retain strong hypertransmission evidence;
    \item reduce, but not eliminate, redundant structural contributions.
\end{itemize}

\end{frame}

\begin{frame}{Initial Grouped Weights}

\textbf{Within the structural group}
\[
S_{\mathrm{structural}}
=
0.45\widetilde V
+
0.275\widetilde C
+
0.275\widetilde H.
\]
\textbf{Within the hypertransmission group}
\[
S_{\mathrm{hyperTD}}
=
0.55\widetilde P
+
0.45\widetilde A.
\]
\textbf{Final group combination}
\[
S
=
0.55S_{\mathrm{structural}}
+
0.45S_{\mathrm{hyperTD}}.
\]

\vspace{0.3cm}

This design reflects the first-scan correlation results:

\begin{itemize}
    \item verticality, continuity, and heterogeneity formed one highly
    correlated structural cluster;
    \item persistence and amplitude formed a second hypertransmission cluster.
\end{itemize}

Grouping reduces the risk of counting closely related structural evidence
three separate times.

\end{frame}

\begin{frame}{Initial Threshold and Spatial Cleanup}

A position is initially marked as barcode-positive when

\[
S(x) > \tau.
\]

The initial proof-of-concept threshold is $\tau=0.0$. This threshold is provisional and has not yet been selected using manual
annotations.

\vspace{0.3cm}

The raw one-dimensional mask is then cleaned using a minimum positive run of 8 positions and maximum gap of 4 positions, along with clipping the first 10 and last 10 positions to account for scanner error


The cleanup order is:

\[
\text{fill short gaps}
\rightarrow
\text{remove short positive runs}.
\]

The final cleaned runs are converted into contiguous horizontal detection
intervals.

\end{frame}

\begin{frame}{First-Run Observations}

The first scan showed encouraging qualitative agreement among multiple
features.

\vspace{0.2cm}

Around the most visually prominent hypertransmission region, all engineered features increased, though how much each increases in EA over barcoding was not investigated.

\vspace{0.3cm}

The strongest observed feature correlations were:
\[
r_{\mathrm{continuity,heterogeneity}}
=
0.920,
\]
\[
r_{\mathrm{verticality,continuity}}
=
0.917,
\]
\[
r_{\mathrm{verticality,heterogeneity}}
=
0.877,
\]
\[
r_{\mathrm{persistence,amplitude}}
=
0.823.
\]
These results support the grouped score design, but they are based on only one
scan and should not yet be treated as general conclusions.

\end{frame}

% ==================================================
% Future work
% ==================================================

\section{Future Work}

\begin{frame}{Future Work}

\begin{enumerate}
    \item \textbf{Expand beyond binary barcoding detection}
    \begin{itemize}
        \item Characterize Normal, Early Atrophy, Barcoding, and possible
        vessel-related structures.
        \item Determine which mathematical features distinguish broad EA from
        narrow vertical barcode patterns.
    \end{itemize}

    \item \textbf{Run the pipeline across all available scans}
    \begin{itemize}
        \item Process all B-scans within each E2E volume.
        \item Compare stability across neighboring scans from the same eye.
        \item Compare center scans and full volumes across subjects.
    \end{itemize}

    \item \textbf{Validate against manual annotations}
    \begin{itemize}
        \item Convert manually labeled intervals into column-level reference
        labels.
        \item Evaluate sensitivity, specificity, precision, F1 score, and
        interval overlap.
    \end{itemize}
\end{enumerate}

\end{frame}

\begin{frame}{Future Weight and Threshold Tuning}

After all available scans are processed, the dataset may contain at least 92*96
scan-level examples, with some volumes containing more than 96 B-scans.

\vspace{0.3cm}

A logistic regression model can then estimate feature weights once more scans have been processed and reliability is assessed.

Training and evaluation must be split by subject or volume rather than by
individual image column to avoid leakage between highly related observations.

Thresholds can then be selected using validation criteria such as:

\begin{itemize}
    \item F1 score;
    \item Youden's index;
    \item a clinically specified sensitivity target.
\end{itemize}

\end{frame}

\begin{frame}{Immediate Next Steps}

\begin{enumerate}
    \item Run the current untuned configuration on additional fast and slow
    progressor volumes.
    \item Inspect whether the five feature signals remain interpretable across
    different scans.
    \item Compare grouped and individual scores.
    \item Quantify within-volume consistency across neighboring B-scans.
    \item Compare automatic intervals with manual Normal, EA, and Barcoding
    labels.
    \item Perform sensitivity analysis for:
    \begin{itemize}
        \item window width and stride;
        \item Gaussian smoothing;
        \item persistence thresholds;
        \item feature weights;
        \item score threshold;
        \item spatial cleanup parameters.
    \end{itemize}
\end{enumerate}

\vspace{0.3cm}

\begin{center}
\textbf{The current results are a proof of concept, not a finalized detector.}
\end{center}

\end{frame}

\end{document}